\section{Fundamentals}






\subsection{Computational Complexity and Algorithmic Approaches}
The factorial scaling of the symmetric TSP, given by $(N-1)!/2$ distinct routes, renders brute-force methods impractical beyond small system sizes. 
Exact algorithms, such as the \textit{Held-Karp dynamic programming approach}, provide solutions but require exponential computational resources, limiting their applicability. 
As a result, most practical implementations rely on \textit{heuristic and approximation algorithms}, which yield near-optimal solutions at significantly reduced computational cost.

Among the most effective heuristic methods are:

\begin{itemize}
    \item \textbf{Greedy algorithms}, which construct solutions iteratively based on local distance minimization.
    \item \textbf{Metaheuristic approaches}, including \textit{simulated annealing, genetic algorithms, and ant colony optimization}, which employ stochastic strategies to escape local minima.
    \item \textbf{Mathematical optimization techniques}, such as \textit{integer linear programming relaxations and cutting-plane methods}, which refine feasible regions iteratively.
\end{itemize}

Although these methods do not guarantee exact solutions, they remain essential in real-world applications where computational constraints necessitate efficient approximations.

\subsection{Physical and Industrial Applications}
Optimization problems analogous to the TSP emerge in numerous domains, particularly in physics and engineering. Examples include:

\begin{itemize}
    \item \textbf{Path optimization in transport networks}, where fuel efficiency and travel time minimization are critical.
    \item \textbf{Manufacturing processes}, where CNC machining and robotic assembly require minimizing tool movement to enhance efficiency.
    \item \textbf{Computational biology}, where the sequencing of DNA fragments or protein folding pathways presents combinatorial structures comparable to TSP-like models.
\end{itemize}

Beyond applied optimization, the TSP is frequently utilized as a \textit{benchmark problem} for evaluating new algorithmic techniques, serving as a standard reference in computational complexity research.

\subsection{Theoretical Significance}
The TSP holds deep implications in computational complexity theory, particularly in relation to the \textit{P vs NP problem}. As an NP-hard problem, the TSP is unlikely to have a polynomial-time solution unless \( P = NP \), a hypothesis that remains one of the most significant open questions in theoretical computer science. A resolution to this problem would not only transform computational mathematics but also have profound consequences for cryptography, artificial intelligence, and algorithmic research.

\subsection{Timeseries Analysis}
\subsubsection{Autocorrelation and effective sample size}
Autocorrelation measures the correlation of a signal with a delayed copy of itself. 

The analysis of autocorrelation is a mathematical tool for identifying repeating patterns or hidden periodicities within a signal obscured by noise.

In the context of Monte Carlo time series, it is used to measure how strongly successive stationary observables $A_n$ are correlated with each other.

The autocorrelation function for a stationary observable is defined as
\begin{equation}
    C_A (t) = \frac{\mathbb{E}[A_n, A_{n+1}] - \mathbb{E}[A]^2}{\sigma_A^2}
    \label{eq:autocorrelation_func}
\end{equation}
The integrated autocorrelation time is convention dependend. Using the lecture convention:
\begin{equation}
    \tau_{\mathrm{int}} = \frac{1}{2} + \sum_{t=1}^{\infty} C_A (t)
    \label{eq:autocorrelation_time}
\end{equation}
For $N$ >> $\tau_{\mathrm{int}}$ the variance of the observables mean $\bar{A}$ can be defined as:
\begin{equation}
    \mathrm{Var}(\bar{A}) = \frac{2 \tau_{\mathrm{int}}}{N} \sigma_A^2
    \label{eq:variance_mean}
\end{equation}

\subsubsection{Blocking}
The ordered series is split into consecutive blocks of length B, each block is replaced by its mean. When $B$ >> $\tau_{\mathrm{int}}$, block means are approximately independent and
the estimated error reaches a plateau. Randomly shuffling first would destroy the time-correlation
information.




